\section{附录：讲座时间线精读}

\subsection{关键时间点与技术主线}
为了便于复习，本节按照时间线提炼讲座主线，并对应到可执行工程动作。此表既可用于复盘，也可直接转为团队内部的学习任务分配。

\begin{longtable}{p{1.8cm}p{3.8cm}p{5.4cm}}
\toprule
\textbf{时间} & \textbf{讲座要点} & \textbf{工程启示} \\
\midrule
\endfirsthead
\toprule
\textbf{时间} & \textbf{讲座要点} & \textbf{工程启示} \\
\midrule
\endhead
00:00--00:03 & 主题定义：Post-Training Verifiable Agents & 明确目标不是 ``更会聊''，而是 ``更会完成可验证任务'' \\
00:03--00:06 & 回顾 RLHF/SFT/PPO/GRPO 基本链路 & 将原有聊天后训练管线扩展到轨迹级学习 \\
00:06--00:10 & 可验证任务需要清晰 verifier & 建立自动验证优先的数据与评估标准 \\
00:10--00:14 & coding agent 环境与工具状态建模 & 数据采集要记录状态转移，不只记录最终答案 \\
00:14--00:18 & verifier 多样性与覆盖挑战 & 使用多维 verifier 组合替代单评分器 \\
00:18--00:22 & 训练与评估分布差异风险 & 评估集要做去污染和难度漂移监控 \\
00:22--00:26 & 泛化不是由单 benchmark 决定 & 建立跨任务、跨工具、跨 harness 的评估矩阵 \\
00:26--00:31 & structured output 与真实可用性问题 & 把格式正确率纳入硬门槛，不靠后处理兜底 \\
00:31--00:34 & holistic evaluation 与 benchmark contamination & 引入持续更新评估池与社区协作审计机制 \\
00:34--00:40 & 进入训练策略讨论前的前提重申 & 先保证数据与评估质量，再谈算法优劣 \\
00:40--00:44 & 两阶段框架：SFT + RL & Light SFT 打底，RL 负责探索与提升 \\
00:44--00:49 & 轻 SFT + 多样样本的重要性 & 避免过重 SFT 把策略空间压窄 \\
00:49--00:54 & 算法仍在演化，理论未收敛 & 保持消融实验与多路线并行探索 \\
00:54--00:58 & train longer / harder tasks / diversity 三观察 & 训练策略需同时考虑步数、难度和多样性 \\
00:58--01:00 & entropy collapse 现象分析 & 建立熵监控与探索保持机制，防止早停滞 \\
01:00--01:03 & on-policy 与 off-policy 对比 & 优先从当前策略轨迹学习，减少反馈错位 \\
01:03--01:05 & clipping 解耦与探索鼓励技巧 & 在稳定与探索之间做可解释的参数平衡 \\
01:05--01:08 & 熵正则等方法缓解塌缩 & 可显式加入熵控制损失，结合任务表现调参 \\
01:08--01:14 & 难题训练集构造与能力分层观察 & 建立按难度分桶的数据调度机制 \\
01:14--01:17 & Q\&A 与系统化建议收束 & 将课程结论转化为团队流程、面板与发布门槛 \\
\bottomrule
\end{longtable}

\subsection{关键术语对照}
\begin{table}[H]
\centering
\begin{tabular}{p{3.0cm}p{2.8cm}p{4.5cm}}
\toprule
\textbf{术语} & \textbf{本讲语境} & \textbf{建议团队内定义} \\
\midrule
Verifiable reward & 可程序化判定的任务反馈 & 可自动复现、可审计、可分解的奖励信号 \\
Environment state & 任务运行上下文全状态 & 含输入、文件、配置、历史动作与观测 \\
Trajectory & 多步行动与反馈序列 & 训练、评估与回放的统一数据主键 \\
Light SFT & 轻量行为先验注入 & 减错而不压缩探索，保持 RL 可提升空间 \\
Entropy collapse & 采样多样性丢失 & 策略退化为单一路径，导致长期停滞 \\
Harness swapping & 执行框架扰动评测 & 测试模型是否依赖偶然实现细节 \\
\bottomrule
\end{tabular}
\caption{课程关键词的工程化对照}
\end{table}

\begin{knowledgebox}{如何使用本附录}
建议将时间线和术语表直接映射到团队周计划：每个时间段对应一个工程行动项，每个术语对应一个内部统一定义，避免跨角色沟通歧义。
\end{knowledgebox}

\subsection{本章小结}
时间线精读的价值在于把 ``听懂'' 变成 ``可执行''。当每个关键观点都对应到具体工程动作，课程内容才能真正进入团队生产体系。

